\section{Discussion}
\label{sec:discussion}

SCOUT is designed for deployments where reconciliation runs as
a continuous background function, most rounds are expected to
be clean, and exact verification is expensive relative to a
fixed-size probe exchange. In this setting, the three probe
outputs serve distinct operational roles: sizing always
provides a severity signal; gating decides whether the current
discrepancy shape supports economical candidate
compression; and triage narrows the verification scope when
the discrepancy is sufficiently concentrated. Table~\ref{tab:probe-outcome-mapping}
summarizes the round-level interpretation. In the common
clean-round case, reconciliation terminates after a single
fixed-cost probe exchange, without invoking downstream
verification. Precise localization and repair within the
candidate scope remain delegated to downstream mechanisms such
as Merkle trees or IBLTs; SCOUT's role ends at diagnosis.


\begin{table}[!htbp]
\caption{Representative operational scenarios and typical probe outcomes.}
\label{tab:probe-outcome-mapping}
\centering
\footnotesize
\setlength{\tabcolsep}{0.2pt}
\begin{tabular}{@{}p{0.26\linewidth}p{0.15\linewidth}p{0.15\linewidth}p{0.45\linewidth}@{}}
\toprule
Scenario & $\hat{D}$ & $g$ & Action \\
\midrule
Normal round
  & $\approx 0$ 
  & -- 
  & {skip downstream verification} \\[2pt]

\shortstack[l]{Shard failure} 
  & varies 
  & \texttt{notgate} 
  & \shortstack[l]{triage, then verify on $\mathcal{C}$} \\[2pt]

\shortstack[l]{One-sided\\backlog}
  & large
  & varies
  & \shortstack[l]{severity alert; triage if\\shape admits, else fallback} \\[2pt]


\shortstack[l]{Diffuse dis-\\crepancy} 
  & \shortstack[l]{moderate\\or large} 
  & \texttt{gate} 
  & \shortstack[l]{bypass triage; use fallback} \\[2pt]

\shortstack[l]{Low obser-\\vability} 
  & varies 
  & \texttt{unknown} 
  & \shortstack[l]{defer, capped triage, or fallback} \\[2pt]

\shortstack[l]{Segmenter\\misconfiguration}
  & $\approx$ total 
  & \texttt{--} 
  & \shortstack[l]{inspect config} \\
\bottomrule
\end{tabular}
\end{table}


Among the probe outcomes in Table~\ref{tab:probe-outcome-mapping}, \texttt{unknown} merits separate discussion. It is an explicit abstention:
  the probe lacks sufficient statistical support for a reliable shape judgment. The most common trigger is a small number of concurrent
faults, which leaves too few sections energized for $K_{\mathrm{eff}}/L$ to reliably separate concentrated from diffuse patterns.
Paradoxically, this is also the regime where triage is most effective, since a small $B$ yields high signal separability and strong candidate
  compression. In deployment, \texttt{unknown} rounds can still be acted upon by combining $\hat{D}$ with other observables such as occupancy:
  deferring when $\hat{D}$ is small, or proceeding with a capped triage pass when $\hat{D}$ is large. What SCOUT withholds is a shape-based
recommendation, not the scale signal itself.

A natural question is whether the discrepancy axes used
in our evaluation---concentration, diffuseness, and side
imbalance---provide a useful abstraction for fault scenarios
encountered in practice. Our results suggest that they capture
several common manifestations relevant to continuous
reconciliation: concentrated discrepancies tend to benefit
from triage, diffuse ones are more likely to be diverted by
gating, and asymmetric ones still admit a stable sizing signal
through de-drifted energy. This abstraction is not intended as
an exhaustive fault taxonomy. One scenario that may appear to
fall outside it is spatially correlated dirty segments in the
underlying key space; however, after hash-based segment
mapping, such correlation often appears as a larger effective
$B$, which is already reflected in the capacity-degradation
dimension characterized in \S\ref{sec:eval-operating}.


Two limitations remain. First, the synthetic and YCSB-driven
workloads do not establish how frequently each discrepancy
class arises in production, nor whether additional classes
emerge under real operational correlations or temporal
patterns. Second, under extreme multi-fault overlap
($B \cdot W$ approaching~$L$), the sparse-overlap
approximation used in the candidate-size analysis begins to
break down, and this regime would benefit from a tighter
characterization.
